%=================================================================
\documentclass[article,accept,moreauthors,pdftex]{Definitions/TESEA} 
\usepackage{bbding}
%=================================================================
\firstpage{1} 
\makeatletter 
%\setcounter{page}{\@firstpage} 
\makeatother
\pubvolume{X}
\issuenum{X}
\articlenumber{XXX}
\pubyear{202X}
\copyrightyear{202X}
\doinum{10.32397/tesea.volX.nX.XXX}
%\externaleditor{Academic Editor: name}
\history{Received: XX Month 202X; Accepted: XX Month 202X; Published: XX Month 202X}
% \updates{yes} % If there is an update available, un-comment this line

%% Uncomment if the new journal that already uses continuous page numbers 
%\continuouspages{yes}

%=================================================================
% Add packages and commands here. The following packages are loaded in our class file: fontenc, calc, indentfirst, fancyhdr, graphicx, lastpage, ifthen, lineno, float, amsmath, setspace, enumitem, mathpazo, booktabs, titlesec, etoolbox, amsthm, hyphenat, natbib, hyperref, footmisc, geometry, caption, url, mdframed, tabto, soul, multirow, microtype, tikz

\usepackage{fancyvrb}
\usepackage{multirow, array} % for tables
\usepackage{float} % to use [H]
\usepackage[utf8]{inputenc}
\usepackage{listings}
\usepackage{amsthm}
%\usepackage{amssymb, amsmath, amsbsy} % symbols
\usepackage{amsbsy} % symbols
\usepackage{upgreek} % for greek letters without italics
\usepackage{mathrsfs} % for font formatting
\usepackage{stackrel} % for \stackbin command
\usepackage{enumitem}
%\usepackage{graphicx}
\usepackage{circuitikz} 
\usepackage{amsfonts}
\usepackage{amssymb}
\usepackage{amsmath}
\usepackage{mathtools}
%\usepackage{url}
%\usepackage{tikz}
\usepackage{pgf}
\usepackage{pgfplots}
\usepgfplotslibrary{groupplots}
\usepackage{xcolor}
\usepackage{pgfplots}
\newcolumntype{C}[1]{>{\centering\arraybackslash}p{#1}}
\usepackage{algorithmic}
\usepackage[lined,boxed,ruled,commentsnumbered]{algorithm2e}
%\usepackage[printwatermark]{xwatermark}
\usetikzlibrary{arrows,shapes.gates.logic.US,shapes.gates.logic.IEC,calc}
%\pgfplotsset{width=7cm,compat=1.8}
\pgfplotsset{compat=1.18} 


\DeclareMathAlphabet{\mathpzc}{OT1}{pzc}{m}{it}

\usetikzlibrary{matrix}
\usepgfplotslibrary{groupplots}
\pgfplotsset{footnotesize,samples=10}

%--------------------
\setitemize{parsep=6pt,itemsep=0pt,leftmargin=*,labelsep=5.5mm}
\setenumerate{parsep=6pt,itemsep=0pt,leftmargin=*,labelsep=5.5mm}
\setlist[description]{itemsep=0mm}  

%=================================================================
%% Please use the following mathematics environments: Theorem, Lemma, Corollary, Proposition, Characterization, Property, Problem, Example, Examples and Definitions, Hypothesis, Remark, Definition, Notation, Assumption
%% For proofs, please use the proof~environment (the amsthm package is loaded by the TESEA class).

%=================================================================
% Full title of~the paper (Capitalized)
\Title{Example article title}


% Author ORCID ID: enter ID or remove command
\newcommand{\orcidauthorA}{0000-0001-9983-XXXX} 
% Add \orcidA{} behind the author's name
% For others ORCIDs define \orcidauthorB, \orcidauthorC, etc.


% Authors, for the paper (add full first names)
\Author{Firstname Lastname$^{1}$, Firstname Lastname$^{2}$ and Firstname Lastname$^{1,*}$\orcidA{}} 
%Please carefully check the accuracy of names and affiliations.



% Authors, for metadata in PDF
\AuthorNames{Lastname, Firstname; Lastname, Firstname; Lastname, Firstname}


% Affiliations / Addresses (Add [1] after \address if there is only one affiliation.)
\address{%
	$^{1}$  \quad Department Name, Institution Name, City, Country.\\
		$^{2}$  \quad Department Name, Institution Name, City, Country. \\
		$^{*}$ \quad Correspondence: e-mail@e-mail.com} 

% Abstract (Do not insert blank lines, i.e., \\) 
\abstract{Please provide an abstract of no more than 300 words. Your abstract should explain the main contributions of your article, and should not contain any material that is not included in the main text. }

% Keywords
\keyword{Keyword1; Keyword2; Keyword3; Keyword4.}

%%% Optional

% \resumen{Proporcione un resumen de no más de 300 palabras. Su resumen debe explicar las principales contribuciones de su artículo y no debe contener ningún material que no esté incluido en el texto principal. El resumen en español es opcional.}
%
% % Palabras claves
% \palabraclave{Despacho económico; Flujo óptimo de potencia; Multiplicadores de Lagrange; Optimización no lineal.}


%%%%%%%%%%%%%%%%%%%%%%%%%%%%%%%%%%%%%%%%%%
\begin{document}
%% How to cite this document
\headrule

\howtocite

%% main text

\section{Introduction}
Your introduction goes here! Some examples of commonly used commands and features are listed below, to help you get started.

\section{Methods and Materials}

Guidelines can be included for standard research article sections, such as this one.

\section{Some \LaTeX{} Examples}
\label{sec:examples}

Use section and subsection commands to organize your document. \LaTeX{} handles all the formatting and numbering automatically. Use ref and label commands for cross-references.

\subsection{Figures and Tables}

Use the table and tabular commands for basic tables --- see Table~\ref{tab:widgets}, for example. You can upload a figure (JPEG, PNG or PDF) using the project menu. To include it in your document, use the includegraphics command as in the code for Figure~\ref{fig:view} below.

\begin{figure}[t]
\centering
\includegraphics[width=0.5\linewidth]{Figure_1.pdf}
\caption{An example figure.}
\label{fig:view}
\end{figure}

\begin{table}[ht]
\centering
\caption{\label{tab:widgets}An example table.}
\begin{tabular}{@{}ll@{}}
\toprule
Item         & Quantity \\ \midrule
Candles      & 4        \\
Fork handles & 3        \\ \bottomrule
\end{tabular}
\end{table}

\subsection{Citations}

LaTeX formats citations and references automatically using the bibliography records in your .bib file, which you can edit via the project menu. Use the cite command for an inline citation, like \cite{Norena2020}. For an inline citation of multiple references, include all in the cite command, like \cite{Norena2020,dhivya2021demand}.

\subsection{Mathematics}

\LaTeX{} is great at typesetting mathematics. Let $X_1, X_2, \ldots, X_n$ be a sequence of independent and identically distributed random variables with $\text{E}[X_i] = \mu$ and $\text{Var}[X_i] = \sigma^2 < \infty$, and let
% 
\begin{equation}
S_n = \frac{X_1 + X_2 + \cdots + X_n}{n}
      = \frac{1}{n}\sum_{i}^{n} X_i \enspace,    
\end{equation}
% 
denote their mean. Then as $n$ approaches infinity, the random variables $\sqrt{n}(S_n - \mu)$ converge in distribution to a normal $\mathcal{N}(0, \sigma^2)$. The text following an equation need not be a new paragraph. Please punctuate equations as regular text.

\subsection{Lists}

You can make lists with automatic numbering \dots

\begin{enumerate}[noitemsep] 
\item Like this,
\item and like this.
\end{enumerate}
\dots or bullet points \dots
\begin{itemize}[noitemsep] 
\item Like this,
\item and like this.
\end{itemize}
\dots or with words and descriptions \dots
\begin{description}
\item[Word] Definition
\item[Concept] Explanation
\item[Idea] Text
\end{description}

\section{Conclusions}
Concisely restate the hypothesis and most important findings. Summarize the major findings, contributions to existing knowledge, and limitations. What are the future directions?


\section*{Acknowledgments}

Additional information can be given in the template, such as to not include funder information in the acknowledgments section.

\vspace{12pt}

\noindent \textbf{Funding:} Please add: “This research received no external funding” or “This research was funded by NAME OF FUNDER grant number XXX.”

\vspace{12pt}

\noindent \textbf{Author contributions:} Each author should have made a substantial contribution to the conception or design of the research; or experimental workl or data analysis and interpretation; or software development; or have drafted the manuscript or significantly revised it. The submitting author is responsible for ensuring that the contributions of all authors are correct. It is expected that all authors will have reviewed, discussed, and agreed to their individual contributions as shared by the submitting author. 

For research articles with several authors, a short paragraph specifying their individual contributions must be provided. The following statements should be used "Conceptualization, A.A. and Y.Y.; Methodology, X.X.; Software, X.X.; Validation, X.X., Y.Y. and Z.Z.; Formal Analysis, X.X.; Investigation, X.X.; Resources, X.X.; Data Curation, X.X.; Writing – Original Draft Preparation, X.X.; Writing – Review \& Editing, X.X.; Visualization, X.X.; Supervision, X.X.; Project Administration, X.X.; Funding Acquisition, Y.Y.”. For more information, please see the CRediT taxonomy website: \url{https://credit.niso.org/}.


\vspace{12pt}

\noindent \textbf{Disclosure statement:} Declare conflicts of interest or state ``The authors declare no conflict of interest.'' Authors must identify and declare any personal circumstances or interest that may be perceived as inappropriately influencing the representation or interpretation of reported research results. Any role of the funders in the design of the study; in the collection, analyses or interpretation of data; in the writing of the manuscript; or in the decision to publish the results must be declared in this section. If there is no role, please state ``The funders had no role in the design of the study; in the collection, analyses, or interpretation of data; in the writing of the manuscript; or in the decision to publish the~results''.
	
%=====================================
% References, variant A: external bibliography
%=====================================
\reftitle{References}
% \externalbibliography{yes}
\bibliography{sample.bib}

	
\end{document}